\chapter{THỰC NGHIỆM VÀ THẢO LUẬN}

\section{Dữ liệu và cách chia tập}

Dữ liệu dùng trực tiếp trong notebook là gói reconia-data trên Kaggle được liên kết trong README của dự án, gồm SROIE, CORD-v2 và WildReceipt. SROIE có các trường company, address, date, total \parencite{huang2019}. CORD cung cấp chú thích văn bản, vị trí và nhóm ngữ nghĩa cho hóa đơn \parencite{park2019}. WildReceipt cung cấp nhãn cho các vùng chữ và được giới thiệu trong công trình SDMG-R \parencite{sun2021}.

\begin{table}[H]
\centering
\renewcommand{\arraystretch}{1.25}
\fontsize{12}{15}\selectfont
\setlength{\tabcolsep}{4pt}
\caption{Số mẫu dùng trong thực nghiệm}
\begin{tabularx}{\textwidth}{>{\raggedright\arraybackslash}X>{\centering\arraybackslash}X>{\centering\arraybackslash}X>{\centering\arraybackslash}X}
\toprule
\textbf{Bộ dữ liệu} & \textbf{Train} & \textbf{Validation} & \textbf{Test} \\
\midrule
SROIE & 501 & 125 & 347 \\
CORD-v2 & 800 & 100 & 100 \\
WildReceipt & 1014 & 253 & 472 \\
Tổng & 2315 & 478 & 919 \\
\bottomrule
\end{tabularx}
\end{table}



CORD-v2 giữ cách chia 800/100/100 có sẵn. SROIE và WildReceipt lấy 20\% tập huấn luyện gốc làm validation với seed 42 và giữ nguyên tập test. Chú thích cung cấp văn bản và tọa độ cho các mô hình sử dụng OCR. Donut và Qwen2.5-VL đọc ảnh và dùng các giá trị tham chiếu để huấn luyện hoặc chấm đầu ra.

\section{Cấu hình huấn luyện}
Các lượt chạy dùng seed 42 và chọn checkpoint có Macro-F1 validation cao nhất. Mỗi cấu hình đặt tối đa 30 epoch. Phiên Qwen trên SROIE dùng thêm ngân sách thời gian để xuất adapter trước giới hạn Kaggle.

\begin{table}[H]
\centering
\renewcommand{\arraystretch}{1.16}
\caption{Cấu hình huấn luyện theo mô hình}
\begin{tabularx}{\textwidth}{X*{3}{>{\centering\arraybackslash}X}}
\toprule
\textbf{Mô hình} & \textbf{Batch} & \textbf{Tích lũy} & \textbf{Tốc độ học} \\
\midrule
RoBERTa & 8 & 1 & \(3 \times 10^{-5}\) \\
BROS & 8 & 1 & \(3 \times 10^{-5}\) \\
PICK & 4 & 2 & \(5 \times 10^{-4}\) \\
LayoutLMv3 & 4 & 2 & \(3 \times 10^{-5}\) \\
Donut & 1 & 8 & \(3 \times 10^{-5}\) \\
Qwen2.5-VL & 1 & 8 & \(10^{-4}\) \\
\bottomrule
\end{tabularx}
\end{table}


\section{Cách tính Precision Recall và F1}
Mô hình được chấm theo cặp trường--giá trị. Giá trị chỉ được tính đúng khi thuộc đúng trường và khớp đáp án sau khi chuyển thành chữ thường và chuẩn hóa khoảng trắng. Mỗi giá trị dự đoán được ghép tối đa với một giá trị tham chiếu.

Gọi TP là số giá trị đúng, FP là số dự đoán thừa hoặc sai và FN là số giá trị bị bỏ sót. Với từng trường, các chỉ số được tính như sau.
\begingroup
\setlength{\jot}{14pt}
\begin{align}
P &= \frac{\mathrm{TP}}{\mathrm{TP}+\mathrm{FP}}, \\
R &= \frac{\mathrm{TP}}{\mathrm{TP}+\mathrm{FN}}, \\
F_1 &= \frac{2\mathrm{TP}}{2\mathrm{TP}+\mathrm{FP}+\mathrm{FN}}.
\end{align}
\endgroup
Nếu mẫu số bằng 0, chỉ số tương ứng bằng 0. Nhóm cộng số đếm theo trường trên toàn bộ tập rồi lấy trung bình các chỉ số giữa những trường xuất hiện trong đáp án hoặc dự đoán. Do đó, Macro-F1 có thể khác F1 tính từ Precision và Recall trung bình. Ví dụ, hai chuỗi 21.60 và 21.6 được tính là khác nhau.

\section{Kết quả trên tập test}
Bảng dùng các lượt chạy mới đã đối chiếu của RoBERTa, BROS, PICK và LayoutLMv3, kết quả Donut trên ba bộ và phiên Qwen2.5-VL đã chấm đủ 347 mẫu SROIE. Các giá trị được làm tròn đến bốn chữ số thập phân. Dấu gạch biểu thị chưa có kết quả test.

\begin{table}[H]
\centering
\renewcommand{\arraystretch}{1.16}
\caption{Kết quả trên tập test SROIE}
\begin{tabularx}{\textwidth}{X*{3}{>{\centering\arraybackslash}X}}
\toprule
\textbf{Mô hình} & \textbf{Precision} & \textbf{Recall} & \textbf{Macro-F1} \\
\midrule
RoBERTa & 0,8759 & 0,8768 & 0,8763 \\
BROS & 0,8774 & 0,8905 & 0,8838 \\
PICK & 0,5567 & 0,6751 & 0,6056 \\
LayoutLMv3 & 0,8828 & 0,8912 & \textbf{0,8869} \\
Donut & 0,7800 & 0,7716 & 0,7758 \\
Qwen2.5-VL & 0,8383 & 0,8379 & 0,8381 \\
\bottomrule
\end{tabularx}
\end{table}

\begin{table}[H]
\centering
\renewcommand{\arraystretch}{1.16}
\caption{Kết quả trên tập test CORD-v2}
\begin{tabularx}{\textwidth}{X*{3}{>{\centering\arraybackslash}X}}
\toprule
\textbf{Mô hình} & \textbf{Precision} & \textbf{Recall} & \textbf{Macro-F1} \\
\midrule
RoBERTa & 0,7975 & 0,8194 & \textbf{0,7984} \\
BROS & 0,7835 & 0,8123 & 0,7870 \\
PICK & 0,3418 & 0,3520 & 0,3362 \\
LayoutLMv3 & 0,7527 & 0,7274 & 0,7344 \\
Donut & 0,5438 & 0,4677 & 0,4859 \\
Qwen2.5-VL & -- & -- & -- \\
\bottomrule
\end{tabularx}
\end{table}

\begin{table}[H]
\centering
\renewcommand{\arraystretch}{1.16}
\caption{Kết quả trên tập test WildReceipt}
\begin{tabularx}{\textwidth}{X*{3}{>{\centering\arraybackslash}X}}
\toprule
\textbf{Mô hình} & \textbf{Precision} & \textbf{Recall} & \textbf{Macro-F1} \\
\midrule
RoBERTa & 0,8602 & 0,8342 & 0,8400 \\
BROS & 0,9224 & 0,8983 & \textbf{0,9094} \\
PICK & 0,5990 & 0,6780 & 0,6313 \\
LayoutLMv3 & 0,9071 & 0,8696 & 0,8728 \\
Donut & 0,5394 & 0,4573 & 0,4861 \\
Qwen2.5-VL & -- & -- & -- \\
\bottomrule
\end{tabularx}
\end{table}

\section{Ví dụ phân tích dự đoán sai}

Ví dụ dưới đây được lấy từ bộ dự đoán theo token đã lưu trong đợt giữa kỳ và minh họa lỗi gán trường của lượt RoBERTa giữa kỳ.

\begin{table}[H]
\centering
\renewcommand{\arraystretch}{1.25}
\fontsize{12}{15}\selectfont
\setlength{\tabcolsep}{4pt}
\caption{Một số lỗi của RoBERTa trên mẫu test CORD-v2 có ID 0}
\begin{tabularx}{\textwidth}{>{\raggedright\arraybackslash}X>{\raggedright\arraybackslash}X>{\raggedright\arraybackslash}X}
\toprule
\textbf{Văn bản} & \textbf{Nhãn GT} & \textbf{Nhãn dự đoán} \\
\midrule
60.000 & B-menu.price & B-menu.unitprice \\
60.000 & B-menu.itemsubtotal & B-menu.price \\
-60.000 & B-sub\_total.discount\_price & B-sub\_total.etc \\
xx7730 & O & B-total.creditcardprice \\
\bottomrule
\end{tabularx}
\end{table}

Ví dụ cho thấy mô hình có thể giữ đúng chữ nhưng gán sai ý nghĩa. Hai đoạn “60.000” là hai vị trí khác nhau trong cùng hóa đơn, mang hai nhãn GT khác nhau. Vì thế, chỉ nhìn danh sách giá trị tiền khó nhận ra nguyên nhân; cần đối chiếu vị trí và nhãn từng đoạn. Một mẫu minh họa không đủ để đại diện cho toàn bộ tập test hoặc chứng minh nguyên nhân chênh lệch giữa các kiến trúc.




\section{Thảo luận và giới hạn}
Trong các lượt hoàn tất, LayoutLMv3 đạt Macro-F1 cao nhất trên SROIE (0,8869), RoBERTa đạt cao nhất trên CORD-v2 (0,7984), và BROS đạt cao nhất trên WildReceipt (0,9094). Qwen2.5-VL đạt Macro-F1 validation 0,8580 tại epoch 8 và test Macro-F1 0,8381. Phiên train dừng theo ngân sách sau 9 epoch hoàn tất. Lượt CORD-v2 ghi nhận loss không hữu hạn ở epoch 1, bước 14. WildReceipt đang chờ lượt chạy.

Đối chiếu CORD-v2 cho thấy một cặp train--test và hai cặp train--validation có văn bản, tọa độ và nhãn trùng nhau. Các ảnh tương ứng gần giống nhau khi kiểm tra theo pixel. Bảng giữ các split đã dùng trong thực nghiệm. Tác động của các mẫu gần trùng lên điểm chưa được đo riêng.

Các lượt được báo cáo dùng một seed. Khởi tạo, đầu vào và cách giải mã khác nhau giữa các mô hình. Với các cấu hình cần OCR, kết quả dựa trên văn bản và tọa độ từ chú thích. Đánh giá trên ảnh mới có OCR tự động là một thiết lập riêng. Với Donut, môi trường của các lượt đã xuất điểm chưa đủ thông tin để xác định ảnh hưởng của cách dịch nhãn trong các phiên bản Transformers. Các điểm trong bảng được đọc từ kết quả đã lưu. Đối chiếu với điểm trong bài báo cần xét cách chấm và bộ dữ liệu của từng nghiên cứu.

\section{Nguồn kết quả thực nghiệm}
Các notebook \href{https://www.kaggle.com/code/vhngthnh/roberta-model}{RoBERTa}, \href{https://www.kaggle.com/code/vhngthnh/bros-model}{BROS}, \href{https://www.kaggle.com/code/vhngthnh/pick-model}{PICK}, \href{https://www.kaggle.com/code/vhngthnh/layoutlmv3-model}{LayoutLMv3} và \href{https://www.kaggle.com/code/vhngthnh/donut-model}{Donut} lưu đầu ra huấn luyện và điểm test. Lượt Qwen SROIE có \href{https://www.kaggle.com/code/phvngtngtm/reconia-qwen-number6-sroie-train}{notebook huấn luyện} và \href{https://www.kaggle.com/code/phvngtngtm/reconia-qwen-number6-sroie-test}{notebook kiểm thử} riêng. Bản thảo chốt số liệu đến ngày 14/11/2026.
